\chapter{Arrays}
\label{ch:arrays}

\begin{goals}
\begin{itemize}
  \item Keeping several values of the same type in one variable
  \item Reaching each element by its number (index) and changing it
  \item Walking through an array with \code{each}
  \item The basic patterns: sum, average, maximum, counting and searching
  \item Working with several parallel arrays
\end{itemize}
\end{goals}

\section{Thirty variables or one array?}

Suppose you want to keep the grades of thirty students and find their
average. With what we have so far, you would have to make thirty variables,
\code{grade1}, \code{grade2}, \ldots{} and add them up thirty times. And if
the class grows to thirty-one? Clearly this road leads nowhere.

The \textbf{array} is the answer to this problem: a single variable that
holds, instead of one value, a list of values of the same type. Instead of
thirty separate boxes, one shelf with thirty numbered compartments.

\section{Making an array and reaching its elements}

An array is made by writing its values between two square brackets,
separated by commas:

\program{Making an array}
\begin{salamcode}
func main:
    grades := [17, 15, 19, 12, 20]
    println grades
    println "First grade:", grades[0]
    println "Third grade:", grades[2]
    println "Number of grades:", grades.len()
end
\end{salamcode}

\begin{salamout}
[17, 15, 19, 12, 20]
First grade: 17
Third grade: 19
Number of grades: 5
\end{salamout}

Several things in this program are new:
\begin{itemize}
  \item \code{println grades} prints the whole array in the same square
        bracket form. It is very useful for a quick look at the contents
        of an array.
  \item \code{grades[0]} is the first element of the array. The number in
        the square brackets is called the \textbf{index}, and it starts
        \emph{from zero}: the first element has index 0, the second has
        index 1, and so on to the end. So \code{grades[2]} is the third
        element.
  \item \code{grades.len()} gives the number of elements in the array. It
        is written with the same dot we saw for \code{to\_int} in chapter
        \ref{ch:types}.
\end{itemize}

\begin{note}[Why from zero?]
Think of the index as ``distance from the start'' rather than ``row
number''. The first element is zero compartments away from the start. Seen
this way, starting at zero is natural, and you saw in chapter
\ref{ch:loops} that the counter of \code{repeat} starts at zero as well;
the two fit together. The last element of an array with 5 elements has
index 4, which is always \code{len() - 1}.
\end{note}

\begin{warn}
If you give an index that is out of range, you get an error. When the index
is a fixed number (for instance \code{grades[5]} for an array of 5
elements, or \code{grades[-1]}), Salam tells you before the program runs:
``index 5 is out of bounds for an array of length 5''. When the index is in
a variable and strays out of range while the program is running, the program
stops on the spot with a runtime error. This is one of the most common
errors in working with arrays, and most of the time it comes from forgetting
``from zero''.
\end{warn}

\section{Changing the elements}

To change one element, we assign to that element through its index. The
array does not need to be \code{mut} for this; \code{mut} is only needed
if you want to replace the whole variable with another array:

\program{Changing an element of an array}
\begin{salamcode}
func main:
    grades := [17, 15, 19, 12, 20]
    println "Before the correction:", grades
    grades[3] = 14
    grades[0] = grades[0] + 1
    println "After the correction:", grades
end
\end{salamcode}

\begin{salamout}
Before the correction: [17, 15, 19, 12, 20]
After the correction: [18, 15, 19, 14, 20]
\end{salamout}

Each element of an array behaves like an ordinary variable: it can be read,
used in arithmetic, and assigned to. The number of elements, however, is
fixed; an array made with five elements has five elements to the end.

\section{Walking through an array with each}

The real power of arrays shows when they are combined with loops. Salam has
a loop made specially for arrays: \kwd{each}. This loop goes round once for
every element of the array, and in each iteration it puts one element into
a variable:

\program{Walking through with each}
\begin{salamcode}
func main:
    fruits := ["apple", "banana", "orange"]
    each fruit in fruits:
        println "-", fruit
    end
end
\end{salamcode}

\begin{salamout}
- apple
- banana
- orange
\end{salamout}

Read \code{each fruit in fruits:} as ``for each fruit in fruits''. In the
first iteration \code{fruit} is ``apple'', in the second ``banana'', and in
the third ``orange''. No index is needed and no counting; Salam hands you
the elements one by one. An array of strings is just like an array of
numbers; the only rule is that all the elements must have the same type.

If you also want the number of each element, write two names inside
parentheses; the first receives the index and the second the element:

\program{Walking through with an index}
\begin{salamcode}
func main:
    fruits := ["apple", "banana", "orange"]
    each (i, fruit) in fruits:
        println i + 1, "-", fruit
    end
end
\end{salamcode}

\begin{salamout}
1 - apple
2 - banana
3 - orange
\end{salamout}

Since the index starts at zero, we added one to it to show a ``human''
number.

\section{The basic patterns of array work}

Most of what we do with arrays comes down to a few simple patterns. All of
them are built from one \code{each} loop and one or two changeable
variables.

\subsection{Sum and average}

\program{The sum and average of the grades}
\begin{salamcode}
func main:
    grades := [17, 15, 19, 12, 20]
    mut total := 0
    each grade in grades:
        total += grade
    end
    count := grades.len()
    average := (total as f64) / (count as f64)
    println "Sum:", total
    println "Average:", average
end
\end{salamcode}

\begin{salamout}
Sum: 83
Average: 16.6
\end{salamout}

This is the adding-up pattern of chapter \ref{ch:loops}, except that this
time the numbers come from an array. If you add ten more grades to the array
tomorrow, not a single other line of the program changes.

\subsection{The maximum}

\program{The hottest day of the week}
\begin{salamcode}
func main:
    temperatures := [21, 25, 19, 30, 27, 24, 22]
    mut highest := temperatures[0]
    mut hottest day := 0
    each (day, temperature) in temperatures:
        if temperature > highest:
            highest = temperature
            hottest day = day
        end
    end
    println "Hottest day: day", hottest day + 1, "at", highest, "degrees"
end
\end{salamcode}

\begin{salamout}
Hottest day: day 4 at 30 degrees
\end{salamout}

This is the ``largest of three numbers'' pattern of chapter
\ref{ch:conditions}: we start with the first element, and every element that
is larger takes its place. This time we also kept the index of the winning
element, so that we know which day it was. To find the minimum, simply turn
\code{>} into \code{<}.

\Needspace{16\baselineskip}
\subsection{Counting}

\program{Counting the passes}
\begin{salamcode}
func main:
    grades := [17, 9, 19, 12, 8, 20, 14]
    mut passed := 0
    each grade in grades:
        if grade >= 10: passed++ end
    end
    println "Passed:", passed, "Failed:", grades.len() - passed
end
\end{salamcode}

\begin{salamout}
Passed: 5 Failed: 2
\end{salamout}

\subsection{Searching}

Is a particular name on the list? For this question we make a boolean
variable that starts as \code{false} and becomes \code{true} the moment the
name is found. Once it is found, we give up the rest of the search with
\code{break}:

\program{Searching a list}
\begin{salamcode}
func main:
    guests := ["Sara", "Reza", "Mina", "Ali"]
    wanted := "Mina"
    mut found := false
    each name in guests:
        if name == wanted:
            found = true
            break
        end
    end
    if found:
        println wanted, "is on the guest list"
    else:
        println wanted, "is not on the guest list"
    end
end
\end{salamcode}

\begin{salamout}
Mina is on the guest list
\end{salamout}

Change the value of \code{wanted} to a name that is not on the list and run
the program again.

\section{An array as a function parameter}

An array can be sent into a function. The type of the parameter is written
as \code{int[:]}, which means ``an array of integers of any length'', and
in the call we write \code{[:]} after the name of the array, meaning ``all
the elements'':

\program{A function that takes an array}
\begin{salamcode}
func sum(numbers: int[:]): int:
    mut total := 0
    each n in numbers:
        total += n
    end
    ret total
end

func main:
    small := [1, 2, 3]
    println sum(small[:])
    purchases := [45, 120, 32]
    println "Purchases total:", sum(purchases[:])
end
\end{salamcode}

\begin{salamout}
6
Purchases total: 197
\end{salamout}

With this function we have written the adding-up pattern once and for all,
and we can use it on any array of any length. (Writing \code{[:]} with
numbers inside the brackets, such as \code{purchases[1:3]}, sends only part
of the array; we do not need that in this book.) The exercises at the end of
the chapter ask you to do the same for the maximum and the average.

\begin{warn}[An array is not copied]
Unlike numbers and strings, an array sent into a function with \code{[:]}
is \emph{not copied}: the function works on the original array itself. If
the function changes one of its elements, the caller's array changes too.
\end{warn}

\section{Several parallel arrays}

Sometimes we have several pieces of data for each thing: a student's name
and grade, a product's name and price. The simple way is to make one array
for each piece of data and agree that element number \code{i} in every
array belongs to the same thing:

\program{Names and grades}
\begin{salamcode}
func main:
    names := ["Sara", "Reza", "Mina"]
    grades := [18, 12, 15]
    each (i, name) in names:
        println name + ":", grades[i]
    end
end
\end{salamcode}

\begin{salamout}
Sara: 18
Reza: 12
Mina: 15
\end{salamout}

The loop goes over the array of names and, with the index it receives, picks
that person's grade out of the second array. This method is simple and
efficient, provided the two arrays always stay the same length and in the
same order.

\begin{deep}[When the count is not known in advance]
The arrays of this chapter have a fixed length: no element can be added to
them or removed. For lists that grow and shrink while the program runs,
Salam has a structure called \code{Vector}, which takes elements with
\code{push} and gives them back with \code{pop}. And to keep related data
together (say a student's name, grade and age), instead of several parallel
arrays you can define a \code{struct}. These are the subject of the next
Salam books; what you learned in this chapter is the foundation of all of
them.
\end{deep}

\section{A more complete program}

The report card of a small class: each person's grade and status, the class
average, and the name of the best student.

\program{A class report card}
\begin{salamcode}
func main:
    names := ["Sara", "Reza", "Mina", "Ali"]
    grades := [18, 12, 15, 19]
    mut total := 0
    mut best := 0

    println "Class report card"
    println "----------------"
    each (i, name) in names:
        grade := grades[i]
        status := grade >= 10 ? "pass" : "fail"
        println name + ":", grade, "(" + status + ")"
        total += grade
        if grade > grades[best]: best = i end
    end
    println "----------------"
    count := grades.len()
    average := (total as f64) / (count as f64)
    println "Class average:", average
    println "Best student:", names[best]
end
\end{salamcode}

\begin{salamout}
Class report card
----------------
Sara: 18 (pass)
Reza: 12 (pass)
Mina: 15 (pass)
Ali: 19 (pass)
----------------
Class average: 16
Best student: Ali
\end{salamout}

This time the variable \code{best} holds the \emph{index} of the best grade
rather than the grade itself; that way, at the end, we can use the same
index to pull the name out of the array of names. Compare this with the
``hottest day'' program, which kept both. Both ways are correct; the choice
depends on what you need at the end.

\begin{summary}
\begin{itemize}
  \item An array is a list of values of the same type:
        \code{grades := [17, 15, 19]}. \code{println} shows the whole
        array and \code{.len()} gives the number of elements.
  \item Indices start at zero: \code{a[0]} is the first element and
        \code{a[a.len() - 1]} the last. An index out of range is an error.
  \item Elements are changed through their index: \code{a[2] = 5}.
        \code{mut} is needed only to replace the whole array.
  \item \code{each element in array:} gives every element in turn, and
        \code{each (index, element) in array:} gives the index too.
  \item The sum, average, maximum, counting and searching patterns are all
        built from one loop and one changeable variable.
  \item An array parameter is written \code{int[:]}, and an array is sent
        into a function as \code{name[:]}.
\end{itemize}
\end{summary}

\section*{Exercises}
\markright{Exercises}

\exercise Make an array of the prices of several items and print the most
expensive and the cheapest, along with the difference between them.

\exercise Write a function called \code{average} that takes an array of
integers and returns their average as a floating-point number. Then write a
function called \code{maximum} using the same pattern.

\exercise Make an array of the temperatures of the seven days of a week and
print how many days were warmer than the week's average. (Hint: you need
two loops; first the average, then the count.)

\exercise Make an array of names and a parallel array of ages. Print the
names of everyone under 18, and at the end say who the oldest person is.

\exercise Make an array of 10 numbers and a new array of 10 zeros. With a
loop, fill the second array with the first one reversed (the last element of
the first in the first element of the second, and so on). Print both.

\exercise What does the program below print? Guess first, then run it:
\begin{salamcode}
func main:
    a := [1, 2, 3, 4, 5]
    each (i, n) in a:
        if (i % 2) == 0: a[i] = n * 10 end
    end
    println a
end
\end{salamcode}
